
\section{Related Work}


\textbf{Multimodal knowledge conflicts.} When information sources disagree, which should a model trust? Prior work has explored this question for text-only language models~\cite{Rich_Knowledge_Sources,xu-etal-2024-knowledge-conflicts,xie2023adaptive}. The problem becomes more pronounced in multimodal settings, where visual perception and textual claims can point to different answers. A number of recent benchmarks document this challenge. HallusionBench~\cite{Hallusionbench}, AutoHallusion~\cite{wu2024autohallusion}, and PhD~\cite{PhD_ChatGPT} report that MLLMs often follow textual context even when it contradicts what is visible in the image. Subsequent analyses suggest this behavior is systematic rather than accidental: models tend to rely on language priors or internal knowledge when facing conflicting evidence~\cite{Insight_Over_Sight,nguyen2025challenges,deng2025words}. More targeted benchmarks explicitly construct multimodal conflicts and measure which modality dominates, further confirming this tendency~\cite{jia2025benchmarking,zhang2025evaluating,zhu2024crossmodalityconflict}. While these studies clearly establish the phenomenon, they mostly describe it at the output level without explaining how conflicts are resolved inside the model. Our work addresses this gap by examining how visual and textual signals interact across layers, showing that many failures occur after visual information has already been encoded correctly.


\textbf{Hallucination mitigation in MLLMs.} Various inference-time methods have been proposed to reduce hallucinations. Contrastive decoding is a common strategy: VCD suppresses tokens preferred under distorted images~\cite{liu2023mitigating}, ICD contrasts standard versus perturbed instructions~\cite{icd2024}, and OPERA penalizes attention over-trust patterns~\cite{huang2024opera}. Recent variants refine these ideas through explicit attention steering~\cite{wang2025ascd}, multi-stage contrast with selective visual inputs~\cite{park2025second}, and probabilistic detection that avoids a second forward pass~\cite{fieback2025ecd}. A separate line of work exploits model depth. DoLa compares predictions from shallow and deep layers to recover factual signals~\cite{dola2024}. DeCo, SHIFT, and related decoders rely on the observation that intermediate layers often preserve stronger visual grounding~\cite{wang2024deco,wang2025shift,layercd_2025}. Training-based approaches reduce hallucinations through preference optimization or contrastive learning but require additional supervision~\cite{wang-etal-2024-mdpo,Yu_2024_CVPR,jiang2024hallucination}.


Most methods correct uniformly, without distinguishing helpful from harmful late-layer processing. Our finding that transition \emph{direction} predicts correctness enables selective intervention only when needed. As we show in Section~\ref{sec:experiments}, uniform correction methods like DoLa can degrade performance when applied indiscriminately.

\textbf{Layer-wise Analysis of Multimodal Models.} Recent studies examine visual information flow in MLLMs, finding that visual signals peak in intermediate layers and weaken in deeper ones~\cite{wang2024valid,knoweledge_evolve2025,yin2025lifting}. Others identify failure modes such as attention drift and modality imbalance that reduce visual grounding~\cite{jiang2025devils,chen2025multimodal}. These findings help explain why models capture visual details early but produce text-driven answers at the end.

Our work extends this line of research. We show that the \emph{direction} of layer-wise transitions predicts final correctness, allowing us to intervene only when late-layer processing undermines visual evidence.
